% Options for packages loaded elsewhere
\PassOptionsToPackage{unicode}{hyperref}
\PassOptionsToPackage{hyphens}{url}
\documentclass[
  ignorenonframetext,
]{beamer}
\newif\ifbibliography
\usepackage{pgfpages}
\setbeamertemplate{caption}[numbered]
\setbeamertemplate{caption label separator}{: }
\setbeamercolor{caption name}{fg=normal text.fg}
\beamertemplatenavigationsymbolsempty
% remove section numbering
\setbeamertemplate{part page}{
  \centering
  \begin{beamercolorbox}[sep=16pt,center]{part title}
    \usebeamerfont{part title}\insertpart\par
  \end{beamercolorbox}
}
\setbeamertemplate{section page}{
  \centering
  \begin{beamercolorbox}[sep=12pt,center]{section title}
    \usebeamerfont{section title}\insertsection\par
  \end{beamercolorbox}
}
\setbeamertemplate{subsection page}{
  \centering
  \begin{beamercolorbox}[sep=8pt,center]{subsection title}
    \usebeamerfont{subsection title}\insertsubsection\par
  \end{beamercolorbox}
}
% Prevent slide breaks in the middle of a paragraph
\widowpenalties 1 10000
\raggedbottom
\AtBeginPart{
  \frame{\partpage}
}
\AtBeginSection{
  \ifbibliography
  \else
    \frame{\sectionpage}
  \fi
}
\AtBeginSubsection{
  \frame{\subsectionpage}
}
\usepackage{iftex}
\ifPDFTeX
  \usepackage[T1]{fontenc}
  \usepackage[utf8]{inputenc}
  \usepackage{textcomp} % provide euro and other symbols
\else % if luatex or xetex
  \usepackage{unicode-math} % this also loads fontspec
  \defaultfontfeatures{Scale=MatchLowercase}
  \defaultfontfeatures[\rmfamily]{Ligatures=TeX,Scale=1}
\fi
\usepackage{lmodern}
\ifPDFTeX\else
  % xetex/luatex font selection
\fi
% Use upquote if available, for straight quotes in verbatim environments
\IfFileExists{upquote.sty}{\usepackage{upquote}}{}
\IfFileExists{microtype.sty}{% use microtype if available
  \usepackage[]{microtype}
  \UseMicrotypeSet[protrusion]{basicmath} % disable protrusion for tt fonts
}{}
\makeatletter
\@ifundefined{KOMAClassName}{% if non-KOMA class
  \IfFileExists{parskip.sty}{%
    \usepackage{parskip}
  }{% else
    \setlength{\parindent}{0pt}
    \setlength{\parskip}{6pt plus 2pt minus 1pt}}
}{% if KOMA class
  \KOMAoptions{parskip=half}}
\makeatother
\setlength{\emergencystretch}{3em} % prevent overfull lines
\providecommand{\tightlist}{%
  \setlength{\itemsep}{0pt}\setlength{\parskip}{0pt}}
\usepackage{bookmark}
\IfFileExists{xurl.sty}{\usepackage{xurl}}{} % add URL line breaks if available
\urlstyle{same}
\hypersetup{
  hidelinks,
  pdfcreator={LaTeX via pandoc}}

\author{\texorpdfstring{}{}}
\date{}

\begin{document}

\begin{frame}[fragile]{Abstract}
\protect\phantomsection\label{abstract}
We quantify how close fixed mathematical constants lie to rationals
\(p/q\) with bounded denominator \(q \leq Q\), and contrast them against
Monte Carlo draws from a \textbf{pooled reference} on \([0,1)\): uniform
samples, fractional parts \(\{\sqrt{k}\}\) for square-free \(k\) from a
fixed list, and (when configured) independent \(\mathrm{Beta}(a,b)\)
draws. The primary statistic is the elementary minimum absolute error
{[} \delta\emph{Q(x) = \min}\{p \in \mathbb{Z},, 1 \leq q \leq Q\}
\left\textbar{} x - \frac{p}{q} \right\textbar, {]} evaluated in
IEEE-754 double precision with a deterministic pseudorandom seed. A
companion scale {[} \mu\emph{Q(x) = \min}\{1 \leq q \leq Q\}~\min\_\{p
\in \mathbb{Z}\} q\^{}2 \left\textbar{} x - \frac{p}{q} \right\textbar{}
{]} records the same best hit weighted by \(q^2\), in the spirit of
Lagrange--Hurwitz quality measures.

Classical Diophantine theory explains asymptotic behaviour as
\(Q \to \infty\): Dirichlet's principle forces
\(\delta_Q(x) \lesssim 1/Q\) for irrationals, Khinchin-type results
describe typical continued-fraction statistics, and badly approximable
numbers (bounded partial quotients, including quadratic irrationals such
as the golden ratio) resist unusually sharp low-height approximations
{[}@khinchin1964continued; @cassels1957{]}. Here we \textbf{fix modest
\(Q\)} and report where named constants sit in the empirical
distribution of \(\delta_Q(X)\) and \(\mu_Q(X)\) relative to the
synthetic baselines---together with optional \textbf{multi-\(Q\)
sensitivity} on the \emph{same} pooled abscissa.

The contribution is pedagogical and computational: a transparent,
test-backed pipeline (\texttt{projects/special\_number\_proximity/src/},
zero-mock tests) links limit theorems to observable finite-\(Q\)
variation. Figures and tabular summaries are produced by
\texttt{scripts/proximity\_monte\_carlo.py} from
\texttt{manuscript/config.yaml}.

\textbf{Keywords:} Diophantine approximation, continued fractions, Monte
Carlo methods, reproducible computing
\end{frame}

\end{document}
